\documentclass{amsart}
% For dealing with references we use the comment environment
\usepackage{verbatim}
\newenvironment{reference}{\comment}{\endcomment}
%\newenvironment{reference}{}{}
\newenvironment{slogan}{\comment}{\endcomment}
\newenvironment{history}{\comment}{\endcomment}

% For commutative diagrams we use Xy-pic
\usepackage[all]{xy}

% We use 2cell for 2-commutative diagrams.
\xyoption{2cell}
\UseAllTwocells

% We use multicol for the list of chapters between chapters
\usepackage{multicol}

% This is generally recommended for better output
\usepackage{lmodern}
\usepackage[T1]{fontenc}

% For cross-file-references

% Package for hypertext links:
\usepackage{hyperref}

% For any local file, say "hello.tex" you want to link to please

% Theorem environments.
%
\theoremstyle{plain}
\newtheorem{theorem}[subsection]{Theorem}
\newtheorem{proposition}[subsection]{Proposition}
\newtheorem{lemma}[subsection]{Lemma}

\theoremstyle{definition}
\newtheorem{definition}[subsection]{Definition}
\newtheorem{example}[subsection]{Example}
\newtheorem{exercise}[subsection]{Exercise}
\newtheorem{situation}[subsection]{Situation}

\theoremstyle{remark}
\newtheorem{remark}[subsection]{Remark}
\newtheorem{remarks}[subsection]{Remarks}

\numberwithin{equation}{subsection}

% Macros
%
\def\lim{\mathop{\mathrm{lim}}\nolimits}
\def\colim{\mathop{\mathrm{colim}}\nolimits}
\def\Spec{\mathop{\mathrm{Spec}}}
\def\Hom{\mathop{\mathrm{Hom}}\nolimits}
\def\Ext{\mathop{\mathrm{Ext}}\nolimits}
\def\SheafHom{\mathop{\mathcal{H}\!\mathit{om}}\nolimits}
\def\SheafExt{\mathop{\mathcal{E}\!\mathit{xt}}\nolimits}
\def\Sch{\mathit{Sch}}
\def\Mor{\mathop{\mathrm{Mor}}\nolimits}
\def\Ob{\mathop{\mathrm{Ob}}\nolimits}
\def\Sh{\mathop{\mathit{Sh}}\nolimits}
\def\NL{\mathop{N\!L}\nolimits}
\def\CH{\mathop{\mathrm{CH}}\nolimits}
\def\proetale{{pro\text{-}\acute{e}tale}}
\def\etale{{\acute{e}tale}}
\def\QCoh{\mathit{QCoh}}
\def\Ker{\mathop{\mathrm{Ker}}}
\def\Im{\mathop{\mathrm{Im}}}
\def\Coker{\mathop{\mathrm{Coker}}}
\def\Coim{\mathop{\mathrm{Coim}}}

% Boxtimes
%
\DeclareMathSymbol{\boxtimes}{\mathbin}{AMSa}{"02}

%
% Macros for moduli stacks/spaces
%
\def\QCohstack{\mathcal{QC}\!\mathit{oh}}
\def\Cohstack{\mathcal{C}\!\mathit{oh}}
\def\Spacesstack{\mathcal{S}\!\mathit{paces}}
\def\Quotfunctor{\mathrm{Quot}}
\def\Hilbfunctor{\mathrm{Hilb}}
\def\Curvesstack{\mathcal{C}\!\mathit{urves}}
\def\Polarizedstack{\mathcal{P}\!\mathit{olarized}}
\def\Complexesstack{\mathcal{C}\!\mathit{omplexes}}
% \Pic is the operator that assigns to X its picard group, usage \Pic(X)
% \Picardstack_{X/B} denotes the Picard stack of X over B
% \Picardfunctor_{X/B} denotes the Picard functor of X over B
\def\Pic{\mathop{\mathrm{Pic}}\nolimits}
\def\Picardstack{\mathcal{P}\!\mathit{ic}}
\def\Picardfunctor{\mathrm{Pic}}
\def\Deformationcategory{\mathcal{D}\!\mathit{ef}}

\setlength{\emergencystretch}{3em}
\begin{document}
\title[Stacks Project, Tag 0C4D]{572 --- Nodes on algebraic curves over arbitrary fields are characterized by quadratic completion, Fitting ideals and geometric branches with delta invariant one}
\maketitle
\noindent Copyright (C) 2005--2025 Johan de Jong. Stacks Project authors. Source: \href{https://stacks.math.columbia.edu/tag/0C4D}{Tag 0C4D}.

\noindent Editorial difficulty note (not author text). Difficulty H2: The equivalence must descend residue-field separability and the Fitting-ideal condition from geometric fibres. The proof then reconciles completion normal forms with normalization lengths and geometric branch counts, including the strict control of delta invariants under field extension..

\noindent Editorial provenance note (not author text). History: excerpt from revision \texttt{a0444\allowbreak 6e57e\allowbreak c1fbc\allowbreak 252a8\allowbreak 71afc\allowbreak ec775\allowbreak 2fb28\allowbreak 07b14}, curves.tex, lines 4173--4301. Reference presentation was adapted; original-author context and core support blocks were retained or added. The mathematical wording is unchanged. Licensed under GNU FDL 1.2 or later; no invariant sections or cover texts. See ../licenses/COPYING. Context blocks reproduce author definitions and supporting results; they are not additional counted problems.

\begin{definition}
\label{definition-nodal}
Let $k$ be a field. Let $X$ be a $1$-dimensional locally algebraic $k$-scheme.
\begin{enumerate}
\item We say a closed point $x \in X$ is a {\it node}, or an
{\it ordinary double point}, or {\it defines a nodal singularity}
if there exists an ordinary double point $\overline{x} \in X_{\overline{k}}$
mapping to $x$.
\item We say the {\it singularities of $X$ are at-worst-nodal} if
all closed points of $X$ are either in the smooth locus of
the structure morphism $X \to \Spec(k)$ or are ordinary double points.
\end{enumerate}
\end{definition}

\begin{lemma}
\label{lemma-nodal-algebraic}
Let $k$ be a field. Let $(A, \mathfrak m, \kappa)$ be a Noetherian
local $k$-algebra. The following are equivalent
\begin{enumerate}
\item $\kappa/k$ is separable, $A$ is reduced,
$\dim_\kappa(\mathfrak m/\mathfrak m^2) = 2$, and there exists a nondegenerate
$q \in \text{Sym}^2_\kappa(\mathfrak m/\mathfrak m^2)$
which maps to zero in $\mathfrak m^2/\mathfrak m^3$,
\item $\kappa/k$ is separable, $\text{depth}(A) = 1$,
$\dim_\kappa(\mathfrak m/\mathfrak m^2) = 2$, and there exists a nondegenerate
$q \in \text{Sym}^2_\kappa(\mathfrak m/\mathfrak m^2)$
which maps to zero in $\mathfrak m^2/\mathfrak m^3$,
\item $\kappa/k$ is separable,
$A^\wedge \cong \kappa[[x, y]]/(ax^2 + bxy + cy^2)$
as a $k$-algebra where $ax^2 + bxy + cy^2$ is a nondegenerate quadratic form
over $\kappa$.
\end{enumerate}
\end{lemma}

\begin{definition}
\label{varieties-definition-delta-invariant}
Let $k$ be a field. Let $X$ be a locally algebraic $k$-scheme.
Let $x \in X$ be a point such that $\mathcal{O}_{X, x}$
is reduced and $\dim(\mathcal{O}_{X, x}) = 1$.
The {\it $\delta$-invariant of $X$ at $x$} is the
$\delta$-invariant of $\mathcal{O}_{X, x}$ as defined in
Definition \ref{varieties-definition-delta-invariant-algebra}.
\end{definition}

\begin{definition}
\label{varieties-definition-delta-invariant-algebra}
Let $A$ be a reduced Nagata local ring of dimension $1$.
The {\it $\delta$-invariant of $A$} is $\text{length}_A(A'/A)$
where $A'$ is as in Lemma \href{https://stacks.math.columbia.edu/tag/0C3S}{lemma pre delta invariant (Tag 0C3S)}.
\end{definition}

\begin{lemma}
\label{lemma-nodal}
Let $k$ be a field. Let $X$ be a $1$-dimensional algebraic $k$-scheme.
Let $x \in X$ be a closed point. The following are equivalent
\begin{enumerate}
\item $x$ is a node,
\item $k \to \mathcal{O}_{X, x}$ is as in Lemma \ref{lemma-nodal-algebraic},
\item any $\overline{x} \in X_{\overline{k}}$ mapping to $x$ defines
a nodal singularity,
\item $\kappa(x)/k$ is separable, $\mathcal{O}_{X, x}$ is reduced, and
the first Fitting ideal of $\Omega_{X/k}$ generates $\mathfrak m_x$
in $\mathcal{O}_{X, x}$,
\item $\kappa(x)/k$ is separable, $\text{depth}(\mathcal{O}_{X, x}) = 1$, and
the first Fitting ideal of $\Omega_{X/k}$ generates $\mathfrak m_x$
in $\mathcal{O}_{X, x}$,
\item $\kappa(x)/k$ is separable and $\mathcal{O}_{X, x}$ is reduced, has
$\delta$-invariant $1$, and has $2$ geometric branches.
\end{enumerate}
\end{lemma}

\begin{proof}
First assume that $k$ is algebraically closed.
In this case the equivalence of (1) and (3) is trivial.
The equivalence of (1) and (3) with (2) holds because the only
nondegenerate quadric in two variables is $xy$ up to change in
coordinates. The equivalence of (1) and (6) is 
Lemma \href{https://stacks.math.columbia.edu/tag/0C1V}{lemma multicross algebra (Tag 0C1V)}.
After replacing $X$ by an affine neighbourhood
of $x$, we may assume there is a closed immersion $X \to \mathbf{A}^n_k$
mapping $x$ to $0$. Let $f_1, \ldots, f_m \in k[x_1, \ldots, x_n]$
be generators for the ideal $I$ of $X$ in $\mathbf{A}^n_k$.
Then $\Omega_{X/k}$ corresponds to the $R = k[x_1, \ldots, x_n]/I$-module
$\Omega_{R/k}$ which has a presentation
$$
R^{\oplus m} \xrightarrow{(\partial f_j/\partial x_i)}
R^{\oplus n} \to \Omega_{R/k} \to 0
$$
(See Algebra, Sections \href{https://stacks.math.columbia.edu/tag/00RM}{section differentials (Tag 00RM)} and
\href{https://stacks.math.columbia.edu/tag/00S0}{section netherlander (Tag 00S0)}.)
The first Fitting ideal of $\Omega_{R/k}$ is thus the ideal
generated by the $(n - 1) \times (n - 1)$-minors of the
matrix $(\partial f_j/\partial x_i)$. Hence (2), (4), (5)
are equivalent by Lemma \href{https://stacks.math.columbia.edu/tag/0C4C}{lemma fitting ideal (Tag 0C4C)} applied
to the completion of $k[x_1, \ldots, x_n] \to R$
at the maximal ideal $(x_1, \ldots, x_n)$.

\medskip\noindent
Now assume $k$ is an arbitrary field.
In cases (2), (4), (5), (6) the residue field $\kappa(x)$ is
separable over $k$. Let us show this holds as well
in cases (1) and (3). Namely, let $Z \subset X$ be the closed subscheme
of $X$ defined by the first Fitting ideal of $\Omega_{X/k}$.
The formation of $Z$ commutes with field extension
(Divisors, Lemma \href{https://stacks.math.columbia.edu/tag/0C3I}{lemma base change and fitting ideal omega (Tag 0C3I)}).
If (1) or (3) is true, then there exists a point
$\overline{x}$ of $X_{\overline{k}}$ such that $\overline{x}$
is an isolated point of multiplicity $1$ of $Z_{\overline{k}}$ (as we have the
equivalence of the conditions of the lemma over $\overline{k}$).
In particular $Z_{\overline{x}}$ is geometrically reduced at $\overline{x}$
(because $\overline{k}$ is algebraically closed). Hence
$Z$ is geometrically reduced at $x$
(Varieties, Lemma \href{https://stacks.math.columbia.edu/tag/0384}{lemma geometrically reduced upstairs (Tag 0384)}).
In particular, $Z$ is reduced at $x$, hence $Z = \Spec(\kappa(x))$
in a neighbourhood of $x$ and $\kappa(x)$ is geometrically reduced
over $k$. This means that $\kappa(x)/k$ is separable
(Algebra, Lemma \href{https://stacks.math.columbia.edu/tag/030W}{lemma characterize separable field extensions (Tag 030W)}).

\medskip\noindent
The argument of the previous paragraph shows that if (1) or (3) holds, then
the first Fitting ideal of $\Omega_{X/k}$ generates $\mathfrak m_x$.
Since $\mathcal{O}_{X, x} \to \mathcal{O}_{X_{\overline{k}}, \overline{x}}$
is flat and since $\mathcal{O}_{X_{\overline{k}}, \overline{x}}$
is reduced and has depth $1$, we see that (4) and (5) hold
(use Algebra, Lemmas \href{https://stacks.math.columbia.edu/tag/033F}{lemma descent reduced (Tag 033F)} and
\href{https://stacks.math.columbia.edu/tag/0337}{lemma apply grothendieck (Tag 0337)}).
Conversely, (4) implies (5) by
Algebra, Lemma \href{https://stacks.math.columbia.edu/tag/031R}{lemma criterion reduced (Tag 031R)}.
If (5) holds, then $Z$ is geometrically reduced at $x$
(because $\kappa(x)/k$ separable and $Z$ is $x$ in a neighbourhood).
Hence $Z_{\overline{k}}$ is reduced at any point $\overline{x}$
of $X_{\overline{k}}$ lying over $x$. In other words, the first
fitting ideal of $\Omega_{X_{\overline{k}}/\overline{k}}$ generates
$\mathfrak m_{\overline{x}}$ in $\mathcal{O}_{X_{\overline{k}, \overline{x}}}$.
Moreover, since
$\mathcal{O}_{X, x} \to \mathcal{O}_{X_{\overline{k}}, \overline{x}}$ is flat
we see that $\text{depth}(\mathcal{O}_{X_{\overline{k}}, \overline{x}}) = 1$
(see reference above).
Hence (5) holds for $\overline{x} \in X_{\overline{k}}$ and we
conclude that (3) holds (because of the equivalence over algebraically
closed fields). In this way we see that (1), (3), (4), (5)
are equivalent.

\medskip\noindent
The equivalence of (2) and (6) follows from
Lemma \href{https://stacks.math.columbia.edu/tag/0C4A}{lemma 2 branches delta 1 (Tag 0C4A)}.

\medskip\noindent
Finally, we prove the equivalence of (2) $=$ (6) with
(1) $=$ (3) $=$ (4) $=$ (5). First we note that the geometric number
of branches of $X$ at $x$ and the geometric number of branches
of $X_{\overline{k}}$ at $\overline{x}$ are equal by
Varieties, Lemma
\href{https://stacks.math.columbia.edu/tag/0C40}{lemma geometric branches and change of fields (Tag 0C40)}.
We conclude from the information available to us at this point
that in all cases this number is equal to $2$.
On the other hand, in case (1) it is clear that $X$ is geometrically
reduced at $x$, and hence
$$
\delta\text{-invariant of }X\text{ at }x \leq
\delta\text{-invariant of }X_{\overline{k}}\text{ at }\overline{x}
$$
by Varieties, Lemma \href{https://stacks.math.columbia.edu/tag/0C3X}{lemma delta invariant and change of fields (Tag 0C3X)}.
Since in case (1) the right hand side is $1$, this
forces the $\delta$-invariant of $X$ at $x$ to be $1$
(because if it were zero, then $\mathcal{O}_{X, x}$ would
be a discrete valuation ring by
Varieties, Lemma \href{https://stacks.math.columbia.edu/tag/0C3U}{lemma delta invariant is zero (Tag 0C3U)}
which is unibranch, a contradiction). Thus (5) holds.
Conversely, if (2) $=$ (5) is true, then assumptions (a), (b), (c) of
Varieties, Lemma \href{https://stacks.math.columbia.edu/tag/0C3P}{lemma geometrically normal in codim 1 (Tag 0C3P)}
hold for $x \in X$ by
Lemma \href{https://stacks.math.columbia.edu/tag/0C4A}{lemma 2 branches delta 1 (Tag 0C4A)}. Thus
Varieties, Lemma
\href{https://stacks.math.columbia.edu/tag/0C3Y}{lemma delta invariant and change of fields better (Tag 0C3Y)}
applies and shows that we have equality in the above displayed inequality.
We conclude that (5) holds for $\overline{x} \in X_{\overline{k}}$
and we are back in case (1) by the equivalence of the conditions
over an algebraically closed field.
\end{proof}
\end{document}
